 \section{Eredmények}
 	
	\subsection{MATLAB implementációk}
	A referenciaként szolgáló MATLAB algoritmus lineáris programszervezést
	alkalmazva az elérhető fajlagos futási idő $\sim100 ms$.
	
	A kód minimális változtatásával elérhető a párhuzamos végrehajtás. Ezt a
	\texttt{for} ciklusok Parallel Toolbox beli \texttt{parfor} utasítására
	cserélve érhetjük el. 4 processzormaggal rendelkező PC esetén ilyenkor
	legjobb esetben a negyedére csökkenhet a futási idő. Valójában ez sose történik meg.
	A MATLAB Parallel Toolbox-a az egyes szimulácókat adott processzormagra osztja el.
	Korábban láttuk, hogy ezen szimulációk lépésigénye nagyban eltérnek egymástól, így előáll az a
	sajnálatos eset, hogy 3 mag tétlenül a negyedikre vár, ami miatt ez nem takinthető járható útnak. 
	
	\subsection{OpenCL implementációk}
	OpenCL keretrendszer segítségével írt programot a GPU-n futtatva a
	\ref{table:openresult}. táblázatban látható eredményeket kapjuk.
	Csupán a globális memóriát használva a referenciához képest romlik a
	teljesítmény. Ezt a videókártya prediktív cache nélküli kialakításának és a
	globális memórájának a kiéheztetésnek tudhatjuk be.
	A lokális memória használata a futási időt drasztikusan le tudja
	csökkenteni, ami a korábban ismertetett memória szervezési megfontolások
	helyességét igazolja.
	
	
	
	
% 	\begin{figure*}[!t]
% 		\centering
% 		\subfloat[Mérési eredmény]{\includegraphics[width=0.45\columnwidth]{kepek/eps/afm200.eps}%
% 		\label{fig_first_case}}
% 		\hfil
% 		\subfloat[Szimulációs eredmény]{\includegraphics[width=0.45\columnwidth]{kepek/eps/efm200.eps}%
% 		\label{fig_second_case}}
% 		\caption{\scriptsize Méréssel kapott felület (balra) és szimulációval kapott töltéstérkép}
% 		\label{fig_sim}
% 	\end{figure*}
	
